% \newpage
\section{Types of Collisions}
\subsection{Definitions}
Collisions are typically classified based on the amount of \textbf{kinetic energy} the system has after the 
collision, in comparison to the amount of kinetic energy the system had before the collision. In other words,
\textbf{how does $E_k'$ with $E_{k}$}?

\subsection{Elastic Collisions}
In an elastic collision, the kinetic energy of the system after the collision is \textbf{equal} to the kinetic
energy of the system before the collision. In mathematics, the equation can be represented by:
\begin{align*}
    E_k' &= E_k
\end{align*}

\begin{remark}
    This does not mean the kinetic energy of the system after the collision is \textbf{equal} to the kinetic 
    energy of the system before the collision (Unlike momentum)
\end{remark}

\subsubsection{Steps of the collision} 
\begin{remark}
    This is on a horizontal frictionless surface, so we can ignore gravitational potential energy.
\end{remark}
    \textbf{Before the collision}: The system's mechanical energy is entirely in the form of kinetic energy. \\

    \textbf{First half of the collision}: The interaction forces cause the objects to start deforming. As the objects deform, 
    they transfer $E_k$ to $E_s$.\\

    \textbf{At the approximate midpoint of the collision}: The deformation of the objects is at a maximum. $E_s$ is maximum and $E_k$ is minimum.\\

\textbf{During the second half of the collision}: The restoring forces do \textit{positive work} on the system, transferring 
elastic potential energy \textbf{back into $E_k$}.\\

\textbf{After the collision}: The system's mechanical energy is now entirely $E_k$. At this time, $E_s = 0$, and all $E_s$ has been transferred to $E_k$. 

\subsubsection{Head-on Collision} 
    \textbf{Before the collision}: The system's mechanical energy is entirely $E_k$.\\

\textbf{During the first half of the collision}: The spring gets compressed. $E_k$ is transformed into $E_s$.\\

\textbf{At the mid-point of the collision}: \begin{itemize}
    \item Spring is at the most compressed point.
    \item The distance between the cars is minimized
    \item $\vec{v_{A}} = \vec{v_{B}}$
    \item $E_k$ is minimized
    \item $E_s$ is maximized
\end{itemize}

\textbf{During the second half of the collision}: \begin{itemize}
    \item $\vec{v_A} < \vec{v_B}$
    \item Distance between the carts is increasing
    \item $E_s$ is being transferred back into $E_k$
\end{itemize}

\textbf{After the collision}:
The system is entirely $E_k$ now.

\begin{remark}
    Elastic collisions \textbf{cannot occur} perfectly between macroscopic objects in real life. At least some of the energy will be lost as thermal energy or sound. 
\end{remark}


\subsection{Inelasic Collision}
So for this collision, $E_k'$ is less than $E_k$.\\

It is impossible for a system to have more kinetic energy after the collision than it had before the collision unless:
\begin{enumerate}
    \item One of the objects had \textbf{stored energy} before the collision, which was transferred into kinetic energy during the collision. 
    \item An \textbf{external} force, such as gravity, is doing positive work on the system during the collision. 
\end{enumerate}

\subsubsection{Completely inelastic collision}
In this collision, the maximum possible amount of kinetic energy is lost as a result of the collision. \\

After the collision, the objects involved will be \textbf{stuck together}.\\

\textbf{Apple and Arrow} is an example of this case. \\

The following condition must be met for a perfectly inelastic collision:
\begin{itemize}
    \item $\vec{v_A}' = \vec{v_B}' = \vec{v}'$
\end{itemize}
